\documentclass[11pt]{article}
\usepackage{amsmath,hyperref}
\title{A free lane is not the same as an available shelf}
\date{}
\begin{document}
\maketitle

\paragraph{The published comparison.}
Adalat, Hamel-De le Court and Belardinelli's \href{https://arxiv.org/html/2606.14130v2}{Contract-Based Compositional Shielding}, Appendix 0.C.2, checks whether an unloaded robot already blocking a currently requested loaded carrier actually yields in its RWARE queue-conflict benchmark. The certified library can change selections between episodes started in covered product states (Theorem 5.2); this does not itself certify replacing a contract in a running episode with outstanding work. The current test executes its original queue protocol evaluator with pinned inputs. It is not a replay of its training or all possible profile obligations, and may not represent its default layout.

\paragraph{Native effects and transfer.}
In the authored layout an outgoing unloaded vehicle can travel through a requested standing shelf; a spare still needs its position to pick it up. A fail-stop imposed just after entering the cell leaves its body there. No loaded requested carrier yet exists, so the published queue-yield proposition remains satisfied even when the original shelf never reaches its destination. A control case places an unloaded blocker in front of an actual loaded requested carrier in the queue: the exact published predicate flags the blocker when it stays and clears the flag when a legal move vacates the cell.

The same stay-clear rule from \href{https://github.com/eternal-november/eos-research-programmes/pull/417}{the fixed-shelf MRWS result} uses a different source interpretation here: the protected source is the actual current mobile shelf corresponding to an unresolved queue entry, and the old robot may otherwise pass underneath it. In the early fault case this rule picks a turn away before forward travel; the fault still occurs at the same tick. The spare picks up the shelf, delivers it at tick 5, and unloads it on another cell at tick 8. No additional slot, scope, shelf, or source of knowledge is added. In a no-fault case the full return takes 15 ticks rather than 10, so the guard is a real tradeoff, not dominance. A separately loaded outgoing robot cannot enter the standing shelf in the native kernel anyway; the added guard makes no difference to its trace. The native movement check itself was exercised in a direct attempted entry control.

A second fail-stop is applied after the spare's first successful delivery. The simulator replaces the request and pays half the two-stage reward. The delivered shelf nonetheless remains attached to the stopped robot at the output cell. Without physical motion and a legal drop, the trailing half of the reward cannot occur. The new queue state cannot substitute for a discharge of the actual return duty. In this trace the queue-yield check still passes; its premise no longer recognizes the old shelf as requested. The closed controller loses the ability to reach the goal cell for further delivery while that body occupies it. The pre-entry guard cannot free a body or unload a shelf after an arbitrary stop.

\paragraph{Scope and evidence.}
The pinned upstream RWARE agent, shelf, grid, highway, and full step method bodies were executed unchanged through syntax-tree extraction. The tested fixtures use exactly one acting vehicle per tick, a checked one-mover graph shim, deterministic one-candidate request replacement, and no Gym observation or learning stack. Faults are authored forced no-ops, not built-in or hardware failures. The published method and its three exact helpers were executed, not a hand-coded approximation of its safety proposition. The test audits 56 paired/control trajectories, in total 840 step calls, plus three isolated controls. Fault ticks are opportunities; in late runs an outgoing actor may have already completed the withdrawal before a trigger, and the comparison does not relabel that arm as a realized fault. The direct planner could protect the same physical cell. No claim of historical novelty, full simulator validation, independent transfer or I1/I2 acceptance follows.
\end{document}
